\eventsection{FUNDAMENTA\c{C}\~AO TE\'ORICA}

A modelagem preditiva em sa\'ude tem avan\c{c}ado com a integra\c{c}\~ao de vari\'aveis ambientais, temporais e territoriais, especialmente em doen\c{c}as sens\'iveis ao clima e \`a sazonalidade. No caso da mal\'aria, estudos recentes demonstram que o uso combinado de aprendizado de m\'aquina e organiza\c{c}\~ao espacial dos dados pode melhorar a previs\~ao de casos na Amaz\^onia Legal \cite{monteiro2025integrating}. Al\'em disso, evid\^encias epidemiol\'ogicas apontam que eventos clim\'aticos extremos e condi\c{c}\~oes meteorol\'ogicas influenciam a transmiss\~ao da doen\c{c}a \cite{arisco2025impact}.

Do ponto de vista metodol\'ogico, a compara\c{c}\~ao entre m\'etodos de \textit{ensemble} \'e particularmente relevante para dados tabulares. O Random Forest representa a estrat\'egia cl\'assica de \textit{bagging}, baseada na combina\c{c}\~ao de \'arvores independentes para reduzir vari\^ancia e aumentar robustez \cite{breiman2001random}. J\'a o \textit{gradient boosting} constr\'oi modelos de forma sequencial, concentrando-se nos erros residuais das etapas anteriores, o que pode gerar maior poder discriminativo em rela\c{c}\~oes n\~ao lineares \cite{friedman2001greedy}. Entre as implementa\c{c}\~oes contempor\^aneas, o XGBoost se destaca pela efici\^encia computacional e regulariza\c{c}\~ao \cite{chen2016xgboost}, enquanto o LightGBM apresenta ganho adicional de velocidade por uso de histogramas e parti\c{c}\~oes otimizadas. Estudos com CatBoost tamb\'em refor\c{c}am o potencial de abordagens de \textit{boosting} quando o tratamento de vari\'aveis categ\'oricas e o controle de vazamento de alvo s\~ao adequadamente considerados \cite{prokhorenkova2019catboostunbiasedboostingcategorical}.

Em problemas de classes raras, a literatura recomenda cautela na escolha das m\'etricas de avalia\c{c}\~ao. \citeonline{5128907} argumentam que distribui\c{c}\~oes desbalanceadas exigem leitura espec\'ifica do comportamento do modelo sobre a classe minorit\'aria. Revis\~oes posteriores e estudos sistem\'aticos mostram que o melhor desempenho depende da intera\c{c}\~ao entre separabilidade dos dados, propor\c{c}\~ao da classe rara e estrat\'egia de reamostragem ou pondera\c{c}\~ao \cite{branco2015surveypredictivemodellingimbalanced,johnson2019survey,BUDA2018249}. Nessa linha, \citeonline{saito2015precision} mostram que a curva Precision-Recall pode ser mais informativa do que a ROC quando a preval\^encia do evento \'e baixa. Complementarmente, \citeonline{davis2006relationship} discutem a rela\c{c}\~ao entre essas curvas e refor\c{c}am a necessidade de interpretar desempenho para al\'em de medidas agregadas simplificadas.

Assim, a fundamenta\c{c}\~ao te\'orica deste estudo combina dois eixos. O primeiro \'e epidemiol\'ogico, ao reconhecer que a mal\'aria no Amazonas responde \`a intera\c{c}\~ao entre din\^amica recente de casos e condi\c{c}\~oes clim\'aticas. O segundo \'e computacional, ao assumir que a compara\c{c}\~ao entre \textit{bagging} e \textit{boosting} deve ser realizada com protocolo temporal rigoroso e com m\'etricas apropriadas para eventos raros. Embora o problema aqui tratado seja bin\'ario, h\'a ainda formula\c{c}\~oes cl\'assicas para extens\~oes de AUC em cen\'arios multiclasse, o que ilustra a centralidade dessa fam\'ilia de m\'etricas na literatura de avalia\c{c}\~ao \cite{article}.
